\section*{Bab \ref{ch:b1:structure-spectroscopy} --- Spektroskopi untuk Menentukan Struktur}

\begin{solution}{exo:b1:structure-spectroscopy:1}
$\varepsilon = A/(\ell c) = 0.64/(1.00 \times \num{2.0e-5}) =
\qty{3.2e4}{L/(mol.cm)}$. Tiga kali konsentrasinya: $A = 1.92$, jika
hukum itu masih berlaku pada \omterm{def:b1:structure-spectroscopy:absorbance}{absorbans} tersebut.
\end{solution}

\begin{solution}{exo:b1:structure-spectroscopy:2}
\ce{C6H6}: $D = (12 + 2 - 6)/2 = 4$, benzena (satu cincin dan tiga ikatan
$\pi$). \ce{C4H8O2}: 1, etil etanoat atau asam butanoat. \ce{C3H7NO}:
$(6 + 2 + 1
- 7)/2 = 1$, propanamida. \ce{C2H3Cl}: 1, kloroetena. \ce{C10H14O}:
$(20 + 2 - 14)/2 = 4$: karvon mempunyai satu cincin, dua C=C, dan satu
C=O.
\end{solution}

\begin{solution}{exo:b1:structure-spectroscopy:3}
Propanon: satu sinyal (6H). Asam etanoat: dua (3H, 1H). 2-Metilpropana:
dua (9H, 1H). 1,2-Dikloroetana: satu (4H).
\end{solution}

\begin{solution}{exo:b1:structure-spectroscopy:4}
$\delta = 1460/400 = \qty{3.65}{ppm}$; pada \qty{90}{MHz} sinyal itu
terletak $3.65 \times 90 = \qty{329}{Hz}$ dari acuan. \omterm{def:b1:structure-spectroscopy:coupling}{Tetapan kopling}
tetap sama dalam hertz, sedangkan selisih pergeseran bertambah dengan
medan: \omterm{def:b1:structure-spectroscopy:coupling}{multiplet-multiplet} yang bertumpang tindih pada medan rendah
terpisah pada medan tinggi.
\end{solution}

\begin{solution}{exo:b1:structure-spectroscopy:5}
\ce{CH3}: triplet (3H); \ce{CH2} tengah: lima tetangga, sekstet
1\,:\,5\,:\,10\,:\,10\,:\,5\,:\,1 (2H); \ce{CH2Br}: triplet (2H), yang
paling kurang terperisai.
\end{solution}

\begin{solution}{exo:b1:structure-spectroscopy:6}
Etanol: O--H (3666) dan tanpa C=O; propanon: C=O (1738) dan tanpa O--H;
asam etanoat: keduanya, O--H (3582) dan C=O (1788).
\end{solution}

\begin{solution}{exo:b1:structure-spectroscopy:7}
Etil etanoat: kuartet (2H) di dekat 4.1, singlet (3H) di dekat 2.0,
triplet (3H) di dekat 1.3 ppm. Metil propanoat: singlet (3H) \ce{OCH3},
yang kurang terperisai karena oksigen (di dekat 3.7 ppm), kuartet (2H)
\ce{CH2C=O} (di dekat 2.3), triplet (3H) di dekat 1.1. Kuartet adalah
sinyal yang paling kurang terperisai pada senyawa pertama dan tidak pada
senyawa kedua: kuartet di dekat 4 ppm berarti \ce{O-CH2CH3}.
\end{solution}

\begin{solution}{exo:b1:structure-spectroscopy:8}
Garis-garisnya berjarak \qty{0.080}{ppm}, yaitu $0.080 \times 90 = \qty{7.2}{Hz}$:
\omterm{def:b1:structure-spectroscopy:coupling}{tetapan koplingnya}. Triplet mempunyai jarak yang sama: \ce{CH2} dan
\ce{CH3} saling berkopling.
\end{solution}

\begin{solution}{exo:b1:structure-spectroscopy:9}
Pita C=O dan satu sinyal: propanon, \ce{CH3COCH3}. Isomer aldehidanya,
propanal \ce{CH3CH2CHO}, memperlihatkan regangan C--H gugus \ce{CHO} dan
protonnya yang sangat kurang terperisai.
\end{solution}

\begin{solution}{exo:b1:structure-spectroscopy:10}
Kumpulan pertama memecah garis menjadi $n_1 + 1$ garis; masing-masing
dipecah lagi oleh kumpulan kedua menjadi $n_2 + 1$: $(n_1 + 1)(n_2 + 1)$
garis. Jika $J_1 =
J_2$, posisi sebuah garis hanya bergantung pada $k_1 + k_2$, yang
mengambil $n_1 + n_2 +
1$ nilai: aturan $n + 1$ dengan $n = n_1 + n_2$. \ce{CH2} tengah
1-bromopropana: $4 \times 3 = 12$ garis secara umum, enam (sebuah
sekstet) jika kedua tetapan itu sama.
\end{solution}

\begin{solution}{exo:b1:structure-spectroscopy:11}
$D = 0$; sebuah alkohol (pita O--H, tanpa C=O). Doblet (6H): dua \ce{CH3}
ekuivalen pada satu CH; \omterm{def:b1:structure-spectroscopy:coupling}{multiplet} (1H): CH itu; doblet (2H): sebuah
\ce{CH2} di sebelah CH itu yang membawa OH; singlet lebar: OH.
2-Metilpropan-1-ol, \ce{(CH3)2CH-CH2-OH}.
\end{solution}

\begin{solution}{exo:b1:structure-spectroscopy:12}
Setiap spesies menyerap secara bebas: $A = A_1 + A_2 = \ell(\varepsilon_1c_1
+ \varepsilon_2c_2)$. Pada dua panjang gelombang, dengan keempat
koefisien diketahui, dua persamaan linear menghasilkan $c_1$ dan $c_2$.
\end{solution}

\begin{solution}{pb:b1:structure-spectroscopy:1}
\textbf{1.} $0.5690 \times 12.0/44.0 = \qty{0.1552}{g}$.
\textbf{2.} $0.2328 \times 2.0/18.0 = \qty{0.02587}{g}$.
\textbf{3.} $0.2500 - 0.1552 - 0.0259 = \qty{0.0689}{g}$.
\textbf{4.} C \qty{62.1}{\percent}, H \qty{10.3}{\percent}, O
\qty{27.6}{\percent}.
\textbf{5.} C $0.1552/12.0 = 0.01293$, H $0.02587$, O $0.0689/16.0 =
0.00431$ mol: rasio $3.00 : 6.00 : 1$, \ce{C3H6O}.
\textbf{6.} $M = \rho RT/p = 3740 \times 8.314 \times 373.15/(1.000 \times 10^5) =
\qty{116}{g/mol}$ (massa jenis dalam \unit{g/m^3}).
\textbf{7.} $116/58 = 2$: \ce{C6H12O2}, $D = (12 + 2 - 12)/2 = 1$.
\textbf{8.} Regangan C=O.
\textbf{9.} O--H mana pun: bukan asam, bukan alkohol.
\textbf{10.} Regangan C--O yang kuat: C--O sebuah ester.
\textbf{11.} Regangan C--H gugus \ce{CH2} dan \ce{CH3}.
\textbf{12.} Satu ikatan $\pi$, dipakai oleh C=O; dengan C--O yang kuat
dan tanpa O--H, sebuah ester. Asam memerlukan O--H; diol tidak
mempunyai C=O. Keton yang membawa gugus eter akan mempunyai C=O-nya di
dekat \qty{1738}{cm^{-1}} propanon, sedangkan nilai teramati
\qty{1754}{cm^{-1}} dekat dengan \qty{1764}{cm^{-1}} etil etanoat; NMR
pada Bagian III menegaskan ester itu.
\textbf{13.} Lima; $2 + 2 + 2 + 3 + 3 = 12$, kedua belas hidrogen.
\textbf{14.} Sebuah \ce{CH2} dengan tiga tetangga (sebuah \ce{CH3}), yang
terikat pada oksigen (kurang terperisai): \ce{-O-CH2-CH3}.
\textbf{15.} \ce{CH3} gugus etil itu, dengan dua tetangga; kedua sinyal
itu saling berkopling, sehingga $J$-nya satu.
\textbf{16.} Sebuah \ce{CH2} dengan dua tetangga, di sebelah C=O.
\textbf{17.} Sebuah \ce{CH2} dengan lima tetangga: di antara sebuah
\ce{CH2} dan sebuah \ce{CH3}.
\textbf{18.} Sebuah \ce{CH3} dengan dua tetangga, jauh dari gugus fungsi.
\textbf{19.} \ce{CH3CH2O-} dan \ce{-CH2CH2CH3}.
\textbf{20.} Proton-proton 4.13 ppm berada pada karbon yang terikat pada
oksigen, proton-proton 2.28 ppm di sebelah C=O: keduanya penarik
elektron.
\textbf{21.} \ce{CH3CH2-O-CO-CH2CH2CH3} atau \ce{CH3CH2-CO-O-CH2CH2CH3}.
Hanya yang pertama menempatkan \ce{CH2} dengan tiga tetangga pada oksigen
(kuartet pada 4.13 ppm).
\textbf{22.} Propil propanoat adalah yang kedua: \ce{OCH2}-nya akan
berupa triplet dan kuartetnya akan terletak di dekat 2.3 ppm.
\textbf{23.} Butil etanoat, \ce{CH3CO-O-C4H9}, akan memperlihatkan singlet
(3H) di dekat 2.0 ppm dan tanpa kuartet.
\textbf{24.} Etil butanoat, \ce{CH3CH2CH2-CO-O-CH2CH3}.
\textbf{25.} IR: C=O 1754, C--O 1182, C--H 2982, tanpa O--H. NMR: \ce{OCH2}
4.13 (q), \ce{CH2CO} 2.28 (t), \ce{CH2} tengah 1.66 (sekstet), \ce{CH3}
ester 1.25 (t), \ce{CH3} rantai 0.95 (t), \omterm{def:b1:structure-spectroscopy:integration}{integrasi} 2, 2, 2, 3, 3.
\textbf{26.} \textbf{\qty{116}{g/mol}}: etil butanoat, \ce{C6H12O2}.
\end{solution}
